\section{Floyd-Warshall Algorithm}
\label{sec:graph-floyd-warshall}

\problemstatement{Given a weighted directed graph (as an
$n \times n$ adjacency matrix, with $\infty$ for missing edges), find
the shortest distance between \textbf{every pair} of nodes.}

\begin{patternbox}
\emph{``Shortest path between every pair of nodes, simultaneously''}
is solved not by running Section~\ref{sec:graph-dijkstra}'s Dijkstra
from every node (which would cost $O(V \cdot (V+E)\log V)$), but by a
direct dynamic programming recurrence over which nodes are
\textbf{allowed as intermediate stops}: progressively allow node $1$,
then nodes $1$ and $2$, then $1,2,3$, and so on, updating every pair's
shortest distance each time a new node becomes available as a
potential stop along the way.
\end{patternbox}

\subsection*{Deriving the recurrence}

Let $\textit{dist}_k[i][j]$ be the shortest distance from $i$ to $j$
using \textbf{only} nodes $1,\ldots,k$ as intermediate stops (besides
$i,j$ themselves). Either the shortest such path doesn't use node $k$
at all (so $\textit{dist}_k[i][j] = \textit{dist}_{k-1}[i][j]$), or it
does use $k$ -- exactly once, as a single stop, splitting the path into
an $i\to k$ segment and a $k\to j$ segment, both of which only need
nodes $1,\ldots,k{-}1$ as their own intermediates (since they don't
revisit $k$):
\[
  \textit{dist}_k[i][j] = \min\big(\textit{dist}_{k-1}[i][j],\
  \textit{dist}_{k-1}[i][k] + \textit{dist}_{k-1}[k][j]\big).
\]
Implemented in place, with $k$ as the \textbf{outermost} loop, this
needs only a single 2D array (each cell is read and written using
already-updated values from the current and prior $k$, which is safe
because $\textit{dist}[i][k]$ and $\textit{dist}[k][j]$ are never
themselves updated using $k$ as an intermediate -- a path from $i$ to
$k$ using $k$ itself as a stop would be nonsensical).

\subsection*{C++ implementation}

\begin{lstlisting}
#include <bits/stdc++.h>
using namespace std;
const int INF = 1e9;

void floydWarshall(vector<vector<int>>& dist) {
    int n = dist.size();

    for (int k = 0; k < n; k++) {
        for (int i = 0; i < n; i++) {
            for (int j = 0; j < n; j++) {
                if (dist[i][k] < INF && dist[k][j] < INF) {
                    dist[i][j] = min(dist[i][j], dist[i][k] + dist[k][j]);
                }
            }
        }
    }
}

int main() {
    int n = 4;
    vector<vector<int>> dist(n, vector<int>(n, INF));
    for (int i = 0; i < n; i++) dist[i][i] = 0;
    dist[0][1] = 5; dist[1][2] = 3; dist[2][3] = 1; dist[0][3] = 100;

    floydWarshall(dist);
    cout << dist[0][3] << endl; // 9  (via 0->1->2->3: 5+3+1)
    return 0;
}
\end{lstlisting}

\subsection*{Dry run}

\dryrun{Take the $4$-node graph above: $0{\to}1(5)$, $1{\to}2(3)$,
$2{\to}3(1)$, $0{\to}3(100)$.}

Initial $\textit{dist}[0][3]=100$ (direct edge). After allowing $k=1$:
$\textit{dist}[0][2] = \min(\infty, \textit{dist}[0][1]+\textit{dist}[1][2]) =
\min(\infty,5+3)=8$. After allowing $k=2$: $\textit{dist}[0][3] =
\min(100, \textit{dist}[0][2]+\textit{dist}[2][3]) = \min(100,8+1)=9$.
After allowing $k=3$: no further improvement possible through $3$
(it has no outgoing edges in this example). Final
$\textit{dist}[0][3]=9$, matching the path $0\to1\to2\to3$.

\begin{complexitybox}
\textbf{Time Complexity.} $O(V^3)$ -- three nested loops over all
nodes.
\textbf{Space Complexity.} $O(V^2)$ for the distance matrix (updated
in place; no extra space needed per layer of $k$).
\end{complexitybox}

\begin{takeawaybox}
Floyd-Warshall is this book's clearest example of an
\textbf{ending-here}-style DP (Section~\ref{sec:graph-shortest-path-intro}'s
opening idea, generalized) applied to a 2D state $(i,j)$ with a third
dimension -- ``which nodes are allowed as intermediates'' -- collapsed
into a single in-place array by processing $k$ in the outer loop, the
same in-place-update trick used by every space-optimized DP table
throughout this book series.
\end{takeawaybox}
